\section{A collection of formal power series}
\label{sec.Phi}

\subsection[The indexing set P\_K]{The indexing set $\boldsymbol{\calP_K}$}
\label{sub.p}

The power series mentioned above are labeled by a finite set $\calP_K$
that coincides with the set of boundary parabolic $\SL_2(\C)$-representations
of $\G:=\pi_1\big(S^3\ssm K\big)$ (or equivalently, of flat connections
on~$S^3\ssm K$ whose restriction to the peripheral subgroup of $\G$ are
parabolic) whenever the latter is finite. For a hyperbolic knot, this
set has three distinguished elements, denoted $\s_0$, $\s_1$ and $\s_2$,
corresponding respectively to the trivial representation, the geometric
representation (given by the natural embedding of~$\G$ into the isometry group
of~$\mathbb H^3$) and the antigeometric representation, which is its
complex conjugate and corresponds to the geometric representation of the
orientation-reversed hyperbolic knot. We denote by $\calP\RED_K=\calP_K\ssm\{\s_0\}$
the reduced set of non-trivial representations (or connections), and often
number the elements of $\calP_K$ as $\s_0,\s_1,\dots,\s_r$, where
$r:=\big|\calP\RED_K\big|$ will be called the \emph{rank} of the knot. (We hope that the
superscript ``red'' will not confuse the reader into thinking that the representations
in $\calP\RED_K$ are reducible; in fact, quite the opposite is true.)
The points of~$\calP\RED_K$ correspond to the solutions (in~$\C$) of a set of polynomial
equations \big(the so-called Neumann--Zagier equations coming from a triangulation
of~$S^3\ssm K$\big) as explained in detail in Section~\ref{sec.perturbation}.
In particular, $\calP\RED_K$ comes equipped with an action of the absolute
Galois group $\text{Gal}\big(\overline{\BQ}/\BQ\big)$, so to every element
$\s \in \calP\RED_K$ is associated a number field~$F_\s$ \big(given either as the
field generated by the coordinates of the solution of the NZ equations or
as the fixed field of the stabilizer of $\s$ in $\text{Gal}\big(\Qbar/\BQ\big)$\big),
called its trace field, together with an embedding, also denoted~$\s$,
of~$F_\s$ into $\Qbar\subset\C$. The field $F_{\s_1}$
coincides with the trace field~$F_K$ of~$K$ as introduced above and
$F_{\s_2}$ is the same field with the complex conjugate embedding into~$\C$.
Two more important invariants of $\s\in\calP_K\RED$ are an element $\xi_\s$
of the Bloch group (or third $K$-group) of~$F_\s$ and a complexified
volume $\V(\s)=\V(K,\s)$ (= ${\rm i}$ times the usual volume plus the Chern--Simons
invariant), obtained as the image of~$\xi_\s$ under the Borel regulator
map or via the dilogarithm, or alternatively its renormalized version
$v(\s)=v(K,\s)=\V(\s)/2\pi {\rm i}$, which for the geometric representation is
the same as the number $v(K)$ occurring in~\eqref{K1}.
(We will usually omit the letter~$K$ in this and all similar notations when
the knot is not varying.) We extend all of these invariants to $\calP_K$
by setting~$F_{\s_0}=\Q$, $\V(s_0)=0$, $\xi_{\s_0}=0$.
In Section~\ref{sec.perturbation} of Part~II,
we will give more details about the set $\calP_K$ and its invariants, and
also say something about the situation when the variety of parabolic
representations contains positive-dimensional components. In
Sections~\ref{sec.Matrix} and~\ref{sec.perturbation}, we will also describe
a large (matrix- rather than vector-valued) collection of power series
associated to~$K$.
